\section{Hexad--Octad Incidence Morphism and Conservation of the Oriented Channels \(90/120\)}

With the dodecad and the alignment of the
\cref{sec:w12-w24-elevacion} fixed, let \(H\) be a hexad in the HMT coordinates,
let \(\ell(H)\subset D\) be its image, and let
\(O_H=\iota_D(\ell(H))\) be its unique lift. Marking a point
and a pair of points of the residual hexad yields
\[
F_6(H)=\ell(H)\times\binom{\ell(H)}2,
\qquad
|F_6|=6\binom62=90.
\]
If the marked point can traverse the octad:
\[
F_8(H)=O_H\times\binom{\ell(H)}2,
\qquad
|F_8|=8\binom62=120.
\]
The inclusion \(\ell(H)\hookrightarrow O_H\) induces the incidence morphism
\[
 \mathcal I_{6\to8}:F_6(H)\hookrightarrow F_8(H).
\]
Normalizing the difference by the chosen
factor \(24\binom62\) yields
\[
\boxed{
\frac{90-120}{24\binom62}
=\frac{6-8}{24}
=-\frac1{12}.}
\]
The equality is a normalized defect of the combinatorial interface \(6\to8\). The factor \(24\binom62\) forms part of the definition of this normalization; the inclusion alone determines the difference \(-30\). The operator \(\mathcal I_{6\to8}\) has marked incidence configurations as its domain; it is neither the transverse extractor of the dodecaphase state nor the family of angular moments.
